% answers_short.tex —— 2016 年·浙江卷·文科 · 简短答案速查
% 由 tools/make_short_answers.py 自 answers.tex 的 \gksolhead/\gkanswer 提取，
% 并与 2016/2016_数学真题_浙江文_答案要点.md 逐题比对后生成，勿手改（改源 MD 后重跑）。

\begin{gknotice}
\textbf{使用说明}：本节只给最终答案便于快速核对，完整推导见解析版。答案为编写者独立推导，\textbf{非官方标准答案}。
\end{gknotice}

\gkpart{选择题}
\begin{center}\begin{tabular}{ccccccccc}
\toprule
题号 & 1 & 2 & 3 & 4 & 5 & 6 & 7 & 8 \\
\midrule
答案 & C & C & D & B & D & A & B & A \\
\bottomrule
\end{tabular}\end{center}

\gkpart{填空题}
\begin{itemize}[leftmargin=2.6em]
  \item[\textbf{9.}] \(S=72\,\mathrm{cm}^{2}\)，\(V=32\,\mathrm{cm}^{3}\)
  \item[\textbf{10.}] 圆心 \((-2,-4)\)，半径 \(5\)
  \item[\textbf{11.}] \(A=\sqrt{2}\)，\(B=1\)
  \item[\textbf{12.}] \(a=-2\)，\(b=1\)
  \item[\textbf{13.}] \((2\sqrt{7},8)\)
  \item[\textbf{14.}] \(\dfrac{\sqrt{6}}{6}\)
  \item[\textbf{15.}] \(\sqrt{7}\)
\end{itemize}

\gkpart{解答题}
\begin{itemize}[leftmargin=2.6em]
  \item[\textbf{16.}] (1) 见解析版；(2) \(\dfrac{22}{27}\)
  \item[\textbf{17.}] (1) \(a_{n}=3^{\,n-1}\)；(2) \(T_{1}=2\)；当 \(n\geqslant 2\) 时 \(T_{n}=\dfrac{3^{\,n}-n^{2}-5n+11}{2}\)
  \item[\textbf{18.}] (1) 见解析版；(2) \(\dfrac{\sqrt{21}}{7}\)
  \item[\textbf{19.}] (1) \(p=2\)；(2) \((-\infty,0)\cup(2,+\infty)\)
  \item[\textbf{20.}] (1) 见解析版；(2) 见解析版
\end{itemize}
